\newcommand{\ThreeHollowSquares}{\tikz[baseline=-.3ex,x=1ex,y=1ex,line width=.09ex]{\draw(0,0)rectangle(.65,.65);\draw(1.3,0)rectangle(1.95,.65);\draw(.65,1)rectangle(1.3,1.65);}}
\begin{center}\bfseries 12.\end{center}
\noindent\hspace*{\parindent}In den D. A. leistet die Theorie der ternären quadratischen Formen aber den ebenso wichtigen Dienst, durch ihre Anwendung auf die Darstellung einer binären Form durch die besondere ternäre Form \(x^2-2x'x''\) erkennen zu lassen, dass jede Klasse des Hauptgeschlechts binärer Formen durch Duplikation entsteht, mithin \(\delta=g\), also auch \(\gamma=\alpha=\frac12\chi\) und somit für die eine Hälfte aller Gesamtcharaktere stets ein dem betreffenden Charakter entsprechendes Geschlecht vorhanden ist. Das Reziprozitätsgesetz gestattet auch, diese Hälfte der Gesamtcharaktere näher zu bestimmen; sie umfasst diejenigen, für welche
\newpage\SourceLeaf{36}\ArithmeticHead{29}{R}{TERNÄRE QUADRATISCHE FORMEN.\quad DER VI. ABSCHNITT DER D. A.}
\noindent \(\left(\frac{D}{n}\right)=1\) ist, wenn \(n\) eine durch eine Form des entsprechenden Geschlechts darstellbare, zu \(2D\) teilerfremde Zahl bedeutet.

\textsc{Gauss} hat sich nicht getäuscht, wenn er (art. 287 Schluss) meinte, diese so einfach lautenden und doch so tief wurzelnden Sätze zu den allerschönsten der ganzen Arithmetik rechnen zu dürfen.

Nachdem \textsc{Gauss} diesen hochbedeutsamen Punkt festgestellt, führte ihn die Theorie der ternären Formen auch zu andern schönen Ergebnissen. Er bestimmte die Anzahl eigentlicher Darstellungen sowohl einer Zahl wie einer binären Form als Summe dreier Quadrate. Während für die Zahlen von den Formen \(8\kappa,\ 8\kappa+4,\ 8\kappa+7\) keine solche möglich ist, lässt jede andere Zahl \(n\) eine Anzahl von Darstellungen zu, welche bemerkenswerter Weise von der Anzahl der Klassen im Hauptgeschlechte derjenigen binären Formen abhängt, deren Determinante \(-n\) ist. So lieferte schon \textsc{Gauss} ein ausgezeichnetes Beispiel viel allgemeinerer Ergebnisse, zu denen nach ihm namentlich H. St. \textsc{Smith} und \textsc{Minkowski} gelangt sind und bei denen der von \textsc{Eisenstein}\footnote{\textsc{Crelles} Journal f. Mathematik \textbf{35} (1847), S. 120; \textbf{41} (1850), S. 151.}\ neu eingeführte Begriff des Masses von Formen oder Darstellungen als bestimmend in den Vordergrund tritt.

In Ea5 des Nachlasses (W. X 1, S. 80) findet sich ein Satz und Bruchstücke zu seinem Beweise, der als Vorläufer zu dem berühmten Satze der D. A. art. 291 betrachtet werden kann. Aus der Darstellung als Summe dreier Quadratzahlen folgerte \textsc{Gauss} dann leicht auch den zuvor noch unbewiesenen \textsc{Fermat}schen Satz, dass jede Zahl als Summe dreier Dreieckszahlen darstellbar sei. Das {\fontspec{times.ttf}ΕΥΡΗΚΑ}, mit welchem \textsc{Gauss} in Nr. 18 T. unter dem 10. Juli 1796 seinen auf diesen Satz bezüglichen Fund hervorhebt, lässt auf den Wert schliessen, den er ihm beigemessen hat, und gestattet die Folgerung, dass, was er gefunden, neu für ihn war. Wenn es die induktive Tatsache des \textsc{Fermat}schen Ausspruchs war, so musste er ihm damals also noch unbekannt sein. Ein direkter, von der Theorie der ternären quadratischen Formen unabhängiger Beweis ist es gewiss nicht gewesen, ihn würde \textsc{Gauss} sicherlich in seinen D. A. mitgeteilt haben. Andererseits zeigt zwar Nr. 17 T. vom 3. Juli 1796, dass er damals wenigstens schon mit der besonderen ternären Form \(x^2+x'^2+x''^2\) sich befasst und für diese gefunden, was später, als er
\newpage\SourceLeaf{37}\ArithmeticHead{30}{L}{BACHMANN, ÜBER GAUSS’ ZAHLENTHEORETISCHE ARBEITEN.}
\noindent die Theorie der ternären quadratischen Formen regelrecht durchführte, in art. 280 allgemein von ihm festgestellt worden ist. Da er aber für die Sätze der artt. 288, 289 über die Darstellung von Zahlen durch jene Form erst viel später (April 1798, Nr. 84 T.) den sicheren Boden fand (das betreffende Resultat ist zu art. 288 erforderlich), bleibt die Annahme eines Induktionsschlusses wahrscheinlich, es müsste denn \textsc{Gauss} insbesondere der Darstellbarkeit der Zahlen \(8n+3\) irgendwie schon gewiss geworden sein und aus ihr den \textsc{Fermat}schen Satz wie in den D. A. gefolgert haben. Vielleicht findet letztere Meinung eine Stütze in der Bemerkung: »Cap. 5. In demonstratione nostra de connexione discerptionum in \ThreeHollowSquares\ et formarum [Linearformen?] ad formas ubi \(\alpha=8n+3\) respiciendum«, welche nebst mehreren andern unter der Überschrift: »Inserenda in opere meo de Residuis« (wohl einer frühen oder erst beabsichtigten Niederschrift seiner D. A.) bei \textsc{Gauss}’ Auszügen aus den Miscell. Taurin. vol. I—IV zu finden ist. Übrigens liest man auch schon in \textsc{Gauss}’ durchschossenem Exemplar von \textsc{Leiste}, \textit{Arithmetik und Algebra} zu S. 68 (W. X 1, S. 78) den Satz »Jede Zahl besteht aus drei Trigonalzahlen« und im Anschluss daran die Angabe, dass jede Zahl von den Formen \(8n+1,\ 8n+3,\ 8n+5\) als Summe von drei, jede Zahl von der Form \(8n+7\) als Summe von vier Quadraten darstellbar sei, mit Angabe zugleich von deren Parität oder Imparität.

Ferner gewährte die allgemeine Theorie der ternären Formen die Herleitung der Bedingungen für die Auflösbarkeit der Gleichung
\[
ax^2+a'x'^2+a''x''^2=0,
\]
die ebenfalls in \textsc{Leiste} zu S. 111 (W. X 1, S. 86) ausgesprochen sind\footnote{Diese Bedingungen waren schon früher von \textsc{Legendre} angegeben worden, siehe dessen \textit{Recherches d’Analyse indéterminée}, 3. article, Histoire de l’Académie de Paris (1785) 1788, Mémoires, S. 507.}, sowie eine Methode zur Auflösung der allgemeinen Gleichung zweiten Grades mit zwei Unbestimmten in rationalen Werten der letztern, während ihre ganzzahlige Auflösung schon an früherer Stelle der D. A. (art. 216 sqq.) mittels Zurückführung auf die Darstellung einer ganzen Zahl durch eine binäre Form geleistet worden war. Die Aufgabe, die allgemeine Gleichung
\[
ax^2+a'x'^2+a''x''^2+2bx'x''+2b'x''x+2b''xx'=0
\]
\newpage\SourceLeaf{38}\ArithmeticHead{31}{R}{TERNÄRE QUADRATISCHE FORMEN.\quad DER VI. ABSCHNITT DER D. A.}
\noindent in ganzen Zahlen zu lösen, findet sich kurz skizziert auch in einem im Juli 1800 begonnenen Notizbüchlein des Nachlasses (Scheda Ae, S. 4, 5, abgedruckt W. X 1, S. 88) behandelt und auf dieselben Prinzipien zurückgeführt, welche \textsc{Gauss} für die Auflösung der genannten einfacheren Gleichung benutzt hat. Es ist dabei bemerkenswert, dass er die Bedingungen für ihre Auflösbarkeit, die aus den für jene einfachere Gleichung geltenden ableitbar und zuerst von H. St. \textsc{Smith}\footnote{Proceed. of the R. Soc. London \textbf{13}, S. 110, Papers I, S. 410.}\ bekannt gemacht worden sind, a. a. O. in völliger Übereinstimmung mit dem Letztern ausgedrückt hat.

Noch sei bemerkt, dass aus einer kurzen Notiz des Nachlasses (Ed 3, abgedruckt W. X 1, S. 92) hervorgeht, dass \textsc{Gauss} auch schon quadratische Formen von der Art, die man \textsc{Hermite}sche nennt, betrachtet hat (siehe a. a. O. S. 94). Dort gibt er die Transformation einer Form
\[
axx'+bx'y+b'xy'+cyy'
\]
mit reellen \(a,c\) und konjugiert komplexen \(b,b'\), für welche er \(\Delta=bb'-ac\) als Determinante benennt, mittels der Substitution
\[
\begin{aligned}
x&=-u,\ y=t+mu\\
x'&=-u',\ y'=t'+m'u',
\end{aligned}
\]
wobei die akzentuierten Grössen die konjugiert imaginären zu den nicht akzentuierten sind. Eine andere Notiz über »duplicierte quadratische Formen«, in welcher auch schon die kubische Form
\[
x^3+ny^3+n^2z^3-3nxyz
\]
auftritt (im Handbuche 18, Bd, vom Oktober 1805, S. 151) deutet durch den dort auftretenden Ausdruck
\[
a+b\sqrt{-3}\mathbin{.}x+c\sqrt{-7}\mathbin{.}y+d\sqrt{21}\mathbin{.}xy
\]
auf eine Beschäftigung mit den biquadratischen Formen hin, welche Normen komplexer, aus zwei Quadratwurzeln gebildeter Zahlen sind. Dies bestätigt eine Notiz in einem der \textsc{Gauss}schen Handexemplare der D. A., nach welcher jede Primzahl von der Form \(8n+1\) als Norm der komplexen Zahl
\newpage\SourceLeaf{39}\ArithmeticHead{32}{L}{BACHMANN, ÜBER GAUSS’ ZAHLENTHEORETISCHE ARBEITEN.}
\[
\alpha+\beta\sqrt{-1}+\gamma\sqrt{-2}+\delta\sqrt{-1}\sqrt{-2}
\]
in der Gestalt
\[
p=(\alpha^2+\beta^2-2\gamma^2-2\delta^2)^2+8(\alpha\gamma+\beta\delta)^2
\]
darstellbar sei.

\textls[100]{Der sechste Abschnitt} der D. A. enthält nur einige Anwendungen des Vorhergehenden teils zur Zerlegung eines Bruchs \(\frac{m}{n}\) in sogenannte Partialbrüche, deren Nenner die in \(n\) enthaltenen Primzahlpotenzen sind, teils zu seiner Entwicklung in einen periodischen Dezimalbruch und zur Bestimmung der Grösse von dessen Periode aus den Exponenten, zu welchen die Zahl 10 bezüglich jener Potenzen als Moduln gehört; ferner Methoden zur Unterscheidung der Primzahlen von den zusammengesetzten Zahlen und im Zusammenhange damit zur Lösung der Kongruenz \(x^2\equiv A\ (\mathrm{mod.}\ M)\), der Gleichung \(mx^2+ny^2=A\) u. s. w., was alles mehr praktisches als theoretisches Interesse erweckt.
